\section{Nearby contours and prediction}\label{li:section}


A second consequence of the correlation estimates is that a normalized
magnetization can be predicted from a nearby contour. This result
supplements ordered transmission when one compares observations on
successive cuts; it is independent of the stopping-band construction
that follows. We first estimate the difference of two line sums and
then use the minimizing property of conditional expectation.

The relative displacement $h=w/s$ is held fixed while $s\to\infty$.
The proof takes the continuum limit of the two-point sums before
estimating their change under displacement. It therefore requires no
discrete derivative estimate for the two-point function.

\begin{proposition}[Differences of nearby line sums]
\label{tr:line-difference}
Let $J_s,J'_s$ be two parallel horizontal or vertical lattice segments of lengths $O(s)$ whose
matching portions differ by a displacement of length at most $Cw$ and
whose unmatched end pieces have total length $O(w)$. The same statement
is allowed for two concentric square contours of radii comparable to
$s$ and radius difference at most $Cw$, by matching their four sides.
Assume the union has a free neighborhood of width $cs$ and is covered
by a fixed number of buffered planar charts. Set
\[
 M_{J_s}=s^{-7/8}\sum_{J_s}\sigma_x,
 \qquad M_{J'_s}=s^{-7/8}\sum_{J'_s}\sigma_x.
\]
There is a constant $C_0$, depending only on the fixed geometry and the
buffer, such that, for each fixed $h=w/s\in(0,h_0]$,
\begin{equation}
\label{tr:line-difference-eq}
 \limsup_{s\to\infty}
       \|M_{J_s}-M_{J'_s}\|_2\leq C_0 h^{3/8}.
\end{equation}
The law may have arbitrary common pins beyond the buffer or be a mixture
of such laws. The same conclusion holds for nearby full periodic rows of period $W$,
normalized by $s^{-7/8}$, when $W/s$ is fixed and positive; its constant
may depend on this ratio.
\end{proposition}

\begin{proof}
The pointwise comparison in Corollary~\ref{fd:buffered-ratio}, applied to
the nonnegative square of a line difference, reduces each buffered chart
to the whole-plane law at the cost of a fixed factor. A finite chart
decomposition and the triangle inequality in $L^2$ will therefore suffice.

In the whole-plane law the mean is zero. By
Lemma~\ref{fd:static-correlations}, normalized two-point sums on
rescaled segments converge to the corresponding integrals of
\[
 K(x-y)=c_\sigma |x-y|^{-1/4}.
\]
If the endpoint data vary with $s$, first pass to a subsequence realizing
the upper limit and then, after a common translation, to a further
subsequence on which their rescaled lengths and relative endpoints
converge. These data range over a compact set, and all the estimates
below are uniform on that set.
Here is the required control at the diagonal. The uniform bound
$\mathbb E(\sigma_x\sigma_y)\leq C(1\vee|x-y|)^{-1/4}$ shows that
parameter pairs whose rescaled distance is at most $r$ contribute at
most $Cr^{3/4}+o(1)$ to the second moment of a normalized segment sum.
The equal-site terms are $O(s^{-3/4})$. Away from the diagonal the
two-point asymptotic is uniform, so ordinary Riemann sums converge.
These observations justify the integral limit for both self and cross
terms, including shifted segments.

For two matching segments, pair their points by a displacement $v$ with
$|v|\leq Ch$. In the limiting squared difference, split pairs of their
one-dimensional parameters at distance $2Ch$. The near part is bounded
by a constant times
\[
 \int_0^{C'h} r^{-1/4}\,dr=O(h^{3/4});
\]
this bound also applies if one of the shifted kernels has its singularity
at a nearby parameter value. For the far part,
$|\nabla K(u)|\leq C|u|^{-5/4}$, and the mean value theorem bounds the
combined kernel difference in absolute value by $Ch|u|^{-5/4}$.
Its integral is at most
\[
 Ch\int_{2Ch}^{C'}r^{-5/4}\,dr=O(h^{3/4}).
\]
Thus the limiting squared norm of the matching-segment difference is
$O(h^{3/4})$. An unmatched end segment of length $O(h)$ has second
moment $O(h^{7/4})$ by the same kernel integral, and hence contributes
$O(h^{7/8})$ in norm. There are only a bounded number of such pieces.
Since $h^{7/8}\leq h^{3/8}$ for $h\leq1$, the triangle inequality proves
\eqref{tr:line-difference-eq}.

For concentric squares, trim corresponding sides to equal lengths,
translate their matching portions in the perpendicular direction, and
apply the segment result to each side and the short remainders.
For full periodic rows, subdivide both rows at the same fixed locations
into matching segments of length a sufficiently small fixed fraction of
$W$. Each pair, together with its buffer, lies in a planar chart; if
necessary decrease $h_0$ by a constant depending on $W/s$. Apply the
preceding comparison and sum the finitely many segment norms. Mixtures
are permitted since the density comparison is uniform before mixing.
\end{proof}

\begin{corollary}[Prediction from a nearby contour]
\label{tr:line-regularity}
In the setting of Proposition~\ref{tr:line-difference}, let $F$ be a
sigma-field containing the spins of $J'_s$, and write
$P_FX=\mathbb E[X\mid F]$. Then
\begin{equation}
\label{tr:line-regularity-eq}
 \limsup_{s\to\infty}\|M_{J_s}-P_FM_{J_s}\|_2
 \leq C_0(w/s)^{3/8}.
\end{equation}
The normalization on the approximating sum remains $s^{-7/8}$ even
when the two contours have different lengths.
\end{corollary}

\begin{proof}
The nearby sum $M_{J'_s}$ is $F$-measurable. The minimizing property of
conditional expectation gives
$\|M_{J_s}-P_FM_{J_s}\|_2\leq\|M_{J_s}-M_{J'_s}\|_2$.
Use Proposition~\ref{tr:line-difference}.
\end{proof}


\begin{remark}[Order of scales]\label{tr:quantifiers}
The transmission exponent, the aspect ratios in the local mass construction,
and the constants for a fixed family of cuts are chosen first. The bounds
are uniform for sufficiently large patch scales. For each of $N$ concentric cuts
at spacings comparable to $s/N$, the local geometry is the same for every
$N$ and $w/s\asymp N^{-1}$. Thus the constant in
\eqref{tr:transmission-lower} is independent of $N$. For the nearby-line
estimate, fix $N$ and then take the mesh sufficiently fine; this threshold
may depend on $N$. A large fixed $N$ may therefore be selected using the
transmission and line constants before taking the mesh limit.
\end{remark}
